% ── SECTION 1: INTRODUCTION (~500 words) ──────────────────────────────────────

\section{Introduction}
\label{sec:intro}

The experimental confirmation of entanglement nonlocality, culminating in the
2022 Nobel Prize in Physics awarded to Aspect, Clauser, and
Zeilinger~\cite{Zeilinger2023Nobel,Aspect1982,CHSH1969}, has definitively
established that quantum mechanics cannot be completed by local hidden
variables~\cite{Bell1964}.
Quantum properties are not intrinsic to isolated systems; they emerge through
relational interaction.
This is one of the most precisely confirmed facts in all of science.

Yet the formal mathematical structure of quantum mechanics has not kept pace
with this empirical verdict.
The dominant interpretations~---~Copenhagen, many-worlds, and QBism~---~either
defer the ontological question (``don't ask what the wave function is, only what
it predicts'') or dissolve it into a proliferation of unobservable branches.
Relational quantum mechanics~\cite{Rovelli1996,Rovelli2021} correctly identifies
the relational character of quantum states~---~quantum properties are always
relative to an observer or physical system, never absolute~---~but has not been
formalized into a mathematical structure that (i)~links coherence, entanglement,
and decoherence dynamics under a single relational concept, and (ii)~derives
testable predictions distinguishable from standard quantum mechanics.

The result is an \emph{interpretive gap}: the empirical facts mandate a
relational ontology, but no formal framework exists to make that ontology
mathematically precise, computationally tractable, and experimentally
discriminating.
This paper closes that gap.

We introduce the \emph{Relational Coherence Theorem} (\RCT), a formal framework
built directly on the relational structure established by Bell's theorem and
relational quantum mechanics.
The \RCT~rests on four propositions: property indefiniteness (quantum properties
are not intrinsic), relational determination (definite properties emerge only
through relational interaction), coherence--affinity correspondence (the degree
of relational coupling is quantified by the Affinity Tensor~$\alpha$), and
conservation of relational potential (quantum information is conserved under
unitary evolution).
Together these propositions formalize the claim that \emph{relational affinity}
--- the magnitude of coherent off-diagonal coupling between a system and its
environment~---~is the fundamental quantity governing both the emergence of
definite properties and the quantum-to-classical transition under
decoherence~\cite{Zurek2003}.

From the \RCT~we derive the \emph{Relational Affinity Gate}~$\URA(\alpha)$, a
continuously parameterized two-qubit unitary whose entanglement capability is
controlled by the Affinity Tensor value~$\alpha \in [0,\pi/2]$.
Unlike fixed entangling gates~(\textsc{cnot}, \textsc{iswap}, \textsc{cz}),
$\URA$~can be tuned to the coherence regime of the physical device, preserving
relational structure under realistic noise.
We characterize~$\URA$ in the Weyl chamber of
$\mathrm{SU}(4)/[\mathrm{SU}(2)\otimes\mathrm{SU}(2)]$~\cite{Zhang2003WeylChamber}
and demonstrate via numerical simulation in Qiskit and Cirq that $\URA$
achieves measurable process-fidelity advantages over fixed gates in
high-decoherence regimes.

This paper presents the mathematical and computational core of Thesis~1.
The philosophical argument~---~that process ontology is formally required (not
merely permitted) by quantum mechanics, and that Whitehead's process philosophy
and the Bah\'{a}'\'{i} writings constitute independent prior instantiations of
the same relational structure~---~is developed in the companion
paper~\cite{Adams2026Philosophy}.

The paper is organized as follows.
Section~\ref{sec:background} reviews Bell's theorem, relational quantum
mechanics, decoherence theory, and relevant entanglement measures.
Section~\ref{sec:rct} formally states and proves the \RCT.
Section~\ref{sec:gate} defines~$\URA$, derives Bell-state production at
$\alpha = \pi/4$, and characterizes the gate in the Weyl chamber.
Section~\ref{sec:sim} presents numerical simulation results under three noise
channels.
Sections~\ref{sec:disc} and~\ref{sec:conc} discuss implications and conclude.
Full proofs of the four propositions appear in Appendix~\ref{app:rct-proof}.
